% Corrected group report, restyled to match the supplied journal-format example.
\documentclass[11pt,a4paper,twoside,reqno]{amsart}
\usepackage[T1]{fontenc}
\usepackage[utf8]{inputenc}
\usepackage[left=28mm,right=28mm,top=25mm,bottom=25mm,headheight=14pt,headsep=7mm]{geometry}
\usepackage{amsmath,amssymb,graphicx,booktabs,tabularx,array}
\usepackage{microtype,needspace}
\usepackage[section]{placeins}
\usepackage[colorlinks=true,urlcolor=blue,linkcolor=black,citecolor=black]{hyperref}
\graphicspath{{images/}}
\setlength{\tabcolsep}{4pt}
\renewcommand{\arraystretch}{1.12}
\urlstyle{same}
\emergencystretch=1.5em
\numberwithin{equation}{section}
\title[Monthly electricity generation in Kazakhstan]{Forecasting Kazakhstan's Total Monthly Electricity Generation Using SARIMA}
\author[Ualikan, Kuantay, Tumenbayev and Sarbassova]{Narlen Ualikan, Dariya Kuantay, Aibek Tumenbayev, Azhar Sarbassova}
\address{Department of Mathematics, Nazarbayev University, Kazakhstan}
\keywords{Time series, electricity generation, Kazakhstan, SARIMA, seasonal forecasting}
\date{}
\hypersetup{pdftitle={Forecasting Kazakhstan's Total Monthly Electricity Generation Using SARIMA},pdfauthor={Narlen Ualikan, Dariya Kuantay, Aibek Tumenbayev, Azhar Sarbassova}}
\begin{document}
\begin{abstract}
This paper compares seasonal naive forecasts, a growth-adjusted seasonal naive benchmark, and SARIMA for Kazakhstan's national monthly electricity generation. The frozen Ember dataset contains 102 months; the main experiment uses 84 observations from 2019 to 2025 after a source audit. Eight low-order specifications are evaluated on 2024 forecasts; SARIMA$(1,0,0)(0,1,1)_{12}$ with a constant is retained for the 2025 test. An unchanged-code rerun reproduces the original test MAEs: @@BASE_MAE@@ TWh for seasonal naive and @@SARIMA_MAE@@ TWh for SARIMA. A follow-up benchmark adds the latest 12-month mean year-on-year change, using only observations through each cutoff. Its MAE is @@GROWTH_MAE@@ TWh and RMSE is @@GROWTH_RMSE@@ TWh, outperforming SARIMA on this test. Thus, SARIMA's @@GAIN@@\% MAE reduction against the original baseline does not establish an advantage over a simple growth adjustment. There are only 12 test observations and four forecast origins. Source construction and observed interval undercoverage limit stronger conclusions.
\end{abstract}
\maketitle
\section{Introduction}

Estimating electricity generation at the national scale provides an application of seasonal forecasting. Short-term forecasts feed into discussions of expected production, but they do not fully determine dispatch, consumption, or infrastructure investments. The objective is to investigate whether SARIMA improves forecasts of monthly electricity generation in Kazakhstan at horizons of one, two, and three months compared with repeating the corresponding months of the previous year.\par

A univariate approach is taken on purpose. The comparison between a transparent seasonal benchmark and a small classical model uses identical forecast origins and horizons. Do short-term dependence and seasonal differencing add information beyond simply repeating last year? A source check and sensitivity analysis help interpret the comparison. The original results showed underprediction during a year of growth. This prompted a follow-up comparison with a fixed growth-adjusted seasonal naive rule; it was not part of the original model selection and is not tuned on 2025.\par

\section{Data description and quality}

Ember's public monthly electricity data [1] were filtered by \textit{Area = Kazakhstan} and \textit{Electricity source = Total generation}. The variable of interest is \textit{Generation (TWh)}, measured in terawatt-hours. One TWh equals one billion kWh. Dates correspond to months. These data reflect the national total, not just Samruk-Energy's plants.\par

\begin{table}[htbp]
\centering
\caption{Audit of the unchanged national series. The current study uses a frozen data version.}
{\fontsize{9.2}{11.3}\selectfont\begin{tabularx}{\linewidth}{@{}>{\hsize=0.64000\hsize\linewidth=\hsize\raggedright\arraybackslash}X>{\hsize=1.36000\hsize\linewidth=\hsize\raggedright\arraybackslash}X@{}}\toprule



Snapshot audit & Verified result \\

\midrule

Coverage & January 2018-June 2026; 102 months \\

Missingness & No missing months, blank targets or duplicate dates \\

Main experiment & January 2019-December 2025; 84 months \\

Snapshot downloaded & 18 November 2026; preserved as a local CSV \\

\bottomrule\end{tabularx}}
\end{table}



\subsection{Source discrepancy and sample choice}

An absence of empty cells does not guarantee accurate measurement. Ember reports 9.780 TWh for each of January, February, and March 2018. Samruk's original report for January [3] shows 10,357.3 million kWh, or 10.3573 TWh. Its first-quarter report gives 29.3392 TWh, very close to Ember's three-month total of 29.340 TWh. This is consistent with quarterly averaging, but it does not prove how the entire file was constructed. The main experiment excludes 2018; Ember's original data stay untouched. Sensitivity to including 2018 is investigated. The partial year of 2026 falls outside the main experiment so validation and testing can proceed with complete calendar years.\par

\subsection{Time-series patterns}

\begin{figure}[htbp]
\centering
\includegraphics[width=\linewidth,height=5.25cm,keepaspectratio]{01_history.png}
\caption{Main study series. Shaded regions mark validation and test years. The historical display includes all study years; model exploration uses only 2019-2023.}
\end{figure}

\begin{figure}[htbp]
\centering
\includegraphics[width=\linewidth,height=5.10cm,keepaspectratio]{02_seasonality.png}
\caption{Monthly patterns in the initial training sample. Higher winter values and lower summer values suggest an annual seasonal component. The plot does not establish a causal weather effect.}
\end{figure}

The training-set lag-12 autocorrelation comes out to 0.681. The ADF test on levels has p = 0.186; the level-stationarity KPSS test has $p\geq 0.10$. At the 5\% level, neither test rejects its respective null hypothesis, so these results are inconclusive about the level series. After seasonal differencing, the ADF p-value drops to 0.0008 and the KPSS p-value is 0.095. Combined with the plots, this supports considering seasonal differencing. The ADF test uses a constant, AIC lag selection, and maximum lag 8. The KPSS test uses a constant and automatic bandwidth. Table-limit p-values are bounds.\par



\section{Methodology}

Consider monthly national generation y\textsubscript{t}. The seasonal naive forecast uses the value from the same month of the previous year. For forecast origin T and horizons h = 1, 2, 3:\par

\[\widehat{y}_{T+h\mid T}=y_{T+h-12}\tag{1}\]

\Needspace{10\baselineskip}
SARIMA(p,d,q)(P,D,Q)\textsubscript{12} combines ordinary (p,d,q) ARIMA components with seasonal (P,D,Q)\textsubscript{12} components [4]. State-space maximum likelihood estimation is used in statsmodels [5], with restrictions ensuring stationary AR and invertible MA components. For the chosen model, define z\textsubscript{t} = y\textsubscript{t} - y\textsubscript{t-12}. Then:\par

\[z_t=c+\phi z_{t-1}+\varepsilon_t+\theta\varepsilon_{t-12}\tag{2}\]

The innovations $\varepsilon_t$ have zero mean and constant variance under the model. The constant allows an average year-on-year change; it is an intercept, not directly the unconditional mean of z\textsubscript{t}. No logarithm, external regressors, or machine-learning model is used. Eight candidates are fixed before test scoring. The smallest unrounded 2024 validation MAE selects the order. For an exact MAE tie, initial-training AICc breaks the tie only when all tied models share the same differencing orders $(d,D)$; otherwise the original candidate enumeration is retained. There were no tied validation MAEs, so this clarification does not alter the selected model. AICc is not the forecast accuracy criterion. All numerical fits converged.\par

\begin{table}[htbp]
\centering
\caption{Chronological design. Each validation and test year contains four non-overlapping three-month blocks.}
{\fontsize{9.2}{11.3}\selectfont\begin{tabularx}{\linewidth}{@{}>{\hsize=1.02000\hsize\linewidth=\hsize\raggedright\arraybackslash}X>{\hsize=1.41000\hsize\linewidth=\hsize\raggedright\arraybackslash}X>{\hsize=0.57000\hsize\linewidth=\hsize\raggedright\arraybackslash}X@{}}\toprule



Stage & Dates & Observations \\

\midrule

Initial training & January 2019-December 2023 & 60 \\

Validation / order selection & January-December 2024 & 12 \\

Held-out test & January-December 2025 & 12 \\

\bottomrule\end{tabularx}}
\end{table}

For January, April, July, and October, models use data only up to the previous month. Parameters are refitted with extended history, but the chosen SARIMA order stays fixed during testing. Forecasts for a three-month period make no use of actual values within that period. Earlier observed periods can, however, be incorporated into later fits. All three models are therefore evaluated on the same target months at horizons 1, 2, and 3. For the same n forecast errors, MAE and RMSE are defined below; both are measured in TWh.\par

\[
\mathrm{MAE}=\frac{1}{n}\sum_{i=1}^{n}|y_i-\widehat y_i|,
\qquad
\mathrm{RMSE}=\sqrt{\frac{1}{n}\sum_{i=1}^{n}(y_i-\widehat y_i)^2}.
\]


\begin{table}[htbp]
\centering
\caption{Candidate comparison, seasonal period 12 throughout. Seasonal naive validation MAE is 0.508 TWh. Minimum unrounded validation MAE selects the model. Exact ties use initial AICc only within a common $(d,D)$; ties spanning differencing orders retain candidate enumeration. No tie occurred. AICc values across differing orders are descriptive only.}
{\fontsize{9.2}{11.3}\selectfont\begin{tabularx}{\linewidth}{@{}>{\hsize=1.64000\hsize\linewidth=\hsize\raggedright\arraybackslash}X>{\hsize=0.60000\hsize\linewidth=\hsize\raggedright\arraybackslash}X>{\hsize=0.88000\hsize\linewidth=\hsize\raggedright\arraybackslash}X>{\hsize=0.88000\hsize\linewidth=\hsize\raggedright\arraybackslash}X@{}}\toprule



SARIMA specification & Constant & Validation MAE & Initial AICc \\

\midrule

(1,0,0)(0,1,1) & Yes & 0.354 & 50.70 \\

(0,0,1)(0,1,1) & Yes & 0.358 & 53.74 \\

(1,0,1)(0,1,1) & Yes & 0.378 & 50.08 \\

(1,0,0)(1,1,0) & Yes & 0.400 & 52.78 \\

(1,1,0)(0,1,1) & No & 0.411 & 53.42 \\

(0,1,1)(0,1,1) & No & 0.429 & 48.13 \\

(1,1,0)(1,1,0) & No & 0.456 & 56.82 \\

(0,1,1)(1,1,0) & No & 0.465 & 50.61 \\

\bottomrule\end{tabularx}}
\end{table}



\subsection{Growth-adjusted seasonal naive: follow-up comparison}
The rule adds an absolute change in TWh, not a percentage growth rate. At cutoff $T$, let
\[
g_T=\frac{1}{12}\sum_{j=0}^{11}(y_{T-j}-y_{T-j-12}),
\qquad \widehat y^{\mathrm{growth}}_{T+h\mid T}=y_{T+h-12}+g_T,
\quad h=1,2,3.\tag{3}
\]
This uses the latest 12 observed year-on-year changes, requiring 24 months of history. The same $g_T$ is added to all three forecasts from a cutoff. It is recomputed at December 2024, March, June and September 2025, with no within-block updating. The 12-month averaging window and additive form are fixed for this follow-up, not selected by trying alternatives on 2025. No intervals are assigned to this point-forecast benchmark. Growth adjustments at the four cutoffs are @@ADJUSTMENTS@@ TWh, respectively.

\subsection{April 2025 example}
Fit through March 2025 and issue April, May and June together. For the growth benchmark, average April 2024--March 2025 year-on-year changes, each compared with its matching month in April 2023--March 2024. This gives $g_T=@@APRIL_ADJ@@$ TWh. For example, April's adjusted forecast is $@@APRIL_BASE@@+@@APRIL_ADJ@@\approx@@APRIL_GROWTH@@$ TWh. April--June actual generation is revealed only for scoring; April's actual value does not update the May or June forecasts. Later July forecasts may use the completed April--June observations.

@@APRIL_TABLE@@

\section{Out-of-sample results}
Before extending the comparison, the unmodified analysis was rerun on the frozen national CSV. It reproduced the original seasonal naive MAE of @@BASE_MAE@@ TWh and SARIMA MAE of @@SARIMA_MAE@@ TWh (no numerical discrepancy at tolerance $10^{-10}$). The selected specification and its 2024 validation MAE of 0.354 TWh are unchanged. The same four blocks and expanding-history refits are used for all three models.

@@SCORES_TABLE@@

\begin{figure}[htbp]
\centering
\includegraphics[width=\linewidth,height=7.0cm,keepaspectratio]{04_test_forecasts.png}
\caption{Actual and predicted generation in 2025. The x-axis has one tick per month; vertical separators mark new three-month forecast blocks starting in April, July and October. Forecasts are issued through December 2024, March, June and September 2025. Shading shows nominal Gaussian SARIMA prediction intervals. Lines join forecasts from different origins for readability.}
\label{fig:test}
\end{figure}

@@HORIZON_TABLE@@

The growth-adjusted benchmark has the lowest pooled MAE (@@GROWTH_MAE@@ TWh) and RMSE (@@GROWTH_RMSE@@ TWh). Its MAE is @@GROWTH_VS_SARIMA@@\% lower than SARIMA's, a difference of @@MAE_DIFF@@ TWh. SARIMA still improves on the original seasonal naive benchmark by @@GAIN@@\%, but that advantage does not persist against the simple growth adjustment. All three pooled mean errors are positive, indicating underprediction on average; the growth adjustment reduces this bias substantially. Every observed month in 2025 exceeds its corresponding 2024 total, which disadvantages the original seasonal naive rule.

The horizon table contains only four forecasts per horizon. The growth benchmark has lower MAE at horizons one and three; SARIMA has slightly lower MAE at horizon two. The pooled ranking therefore does not imply that one method wins at every horizon or every origin. This is a follow-up comparison motivated by the original test results, not a new untouched confirmatory test. The improvement cannot establish statistical significance or general superiority; another test period would be needed for an independent assessment.

\section{Model estimates and diagnostics}

The estimates used for the diagnostic plot do not use observations from the test year. Table~\ref{tab:coef} reports maximum-likelihood estimates and their asymptotic standard errors. These are not the estimates from each quarterly refit, but from a model fitted through December 2024. The seasonal MA estimate is relatively uncertain, which is consistent with the small number of annual cycles.\par

\begin{table}[htbp]
\centering
\caption{Pre-test fit using January 2019-December 2024. In equation (2), c is the intercept of z, not its unconditional mean.}\label{tab:coef}
{\fontsize{9.2}{11.3}\selectfont\begin{tabularx}{\linewidth}{@{}>{\hsize=1.56000\hsize\linewidth=\hsize\raggedright\arraybackslash}X>{\hsize=0.72000\hsize\linewidth=\hsize\raggedright\arraybackslash}X>{\hsize=0.72000\hsize\linewidth=\hsize\raggedright\arraybackslash}X@{}}\toprule



Parameter & Estimate & Standard error \\

\midrule

Constant c & 0.1214 & 0.0313 \\

AR coefficient phi & 0.3138 & 0.1409 \\

Seasonal MA theta & -0.8640 & 0.3788 \\

Innovation variance (TWh squared) & 0.1162 & 0.0328 \\

\bottomrule\end{tabularx}}
\end{table}

\begin{figure}[htbp]
\centering
\includegraphics[width=\linewidth,height=7.70cm,keepaspectratio]{05_residuals.png}
\caption{Residual sequence and autocorrelation from the pre-test fit. The first 13 residuals are omitted to reduce initialization effects. Remaining residual mean = 0.016 TWh. ACF bounds are approximate individual-lag reference bounds.}
\end{figure}

Ljung-Box statistics are 8.075 at lag 12 (p = 0.621) and 18.888 at lag 24 (p = 0.652), with degrees of freedom adjusted for two AR/MA coefficients. These tests fail to reject the hypothesis of no residual autocorrelation at these lags. They do not establish independence, normality, constant variance, or correct model specification. The residual plot and source-data limitations still apply.\par

The SARIMA 80\% intervals cover 9 of 12 test observations (75.0\%); its 95\% intervals cover 10 of 12 (83.3\%). These are nominal Gaussian prediction intervals conditional on the fitted SARIMA model and estimated parameters. They do not fully account for parameter uncertainty or data revisions. Observed coverage falls below the nominal levels even though point-forecast errors improve. Parameter uncertainty, model misspecification, and data limitations may contribute, but this small test cannot establish the cause. Twelve observations are insufficient to estimate long-run coverage precisely.\par



\section{Sensitivity, limitations and conclusion}

Including the full year 2018 in the training history, while retaining the selected SARIMA specification, gives test MAE of 0.360 TWh and RMSE of 0.441 TWh. Seasonal naive forecasts do not change because their corresponding months are in 2024. The growth-adjusted forecasts also remain unchanged: at the first cutoff their adjustment uses 2023--2024 only, and later cutoffs use more recent observations. SARIMA still beats the original seasonal naive baseline, but remains behind the growth-adjusted benchmark (MAE @@GROWTH_MAE@@ TWh). This is a sensitivity analysis, not a step in model selection.\par

The analysis rests on a frozen data snapshot, not the data vintages available at each historical forecast origin. Later data revisions or source construction may affect published series values. Other repeated observations exist after 2018; only a small portion of the source chain was independently verified. These forecasts should be viewed as forecasts of the published Ember series, not independently validated measurements of physical generation.\par

The study is limited to a single held-out year and four forecast origins. The univariate model does not explicitly account for weather, plant maintenance, fuel availability, trade, or policy changes. Published monthly totals are affected by the varying number of days in each month. Generation is not demand; historical forecasting ability does not establish a causal relationship or an operational planning rule.\par

As an example, refitting the model through December 2025 yields the following three-month forecasts. They use only information available through the study endpoint and are explicitly dated to that forecast origin. They are not live forecasts issued in November 2026 and are not part of model selection or primary testing.\par

\begin{table}[htbp]
\centering
\caption{Dated forecast illustration, conditional on observations through December 2025.}
{\fontsize{9.2}{11.3}\selectfont\begin{tabularx}{\linewidth}{@{}>{\hsize=1.00000\hsize\linewidth=\hsize\raggedright\arraybackslash}X>{\hsize=1.00000\hsize\linewidth=\hsize\raggedright\arraybackslash}X>{\hsize=1.00000\hsize\linewidth=\hsize\raggedright\arraybackslash}X>{\hsize=1.00000\hsize\linewidth=\hsize\raggedright\arraybackslash}X@{}}\toprule



Forecast month & Forecast (TWh) & 95\% lower & 95\% upper \\

\midrule

2026-01 & 11.860 & 11.130 & 12.590 \\

2026-02 & 10.812 & 10.052 & 11.572 \\

2026-03 & 11.367 & 10.605 & 12.129 \\

\bottomrule\end{tabularx}}
\end{table}

\textbf{Conclusion.} SARIMA improves on the original seasonal naive model in the 2025 historical test, with a @@GAIN@@\% reduction in pooled MAE. However, it does not outperform the growth-adjusted seasonal naive benchmark: their MAEs are @@SARIMA_MAE@@ and @@GROWTH_MAE@@ TWh, respectively. The more demanding comparison shows that a simple adjustment for recent annual change is sufficient to outperform the selected SARIMA on this particular year. This benchmark was added after seeing the original results and requires independent evaluation in future work. Residual tests do not detect remaining autocorrelation at the tested lags, but the SARIMA 80\% and 95\% intervals cover only 9/12 and 10/12 test observations. A single test year, four origins and the source audit limit stronger reliability claims. More forecast origins and a more extensive audit of the monthly source series would strengthen the evidence.

\begin{thebibliography}{9}

\bibitem{ref1} Ember. Monthly Electricity Data. Snapshot downloaded 18 November 2026. \newline\url{https://ember-energy.org/data/monthly-electricity-data/}

\bibitem{ref2} Ember. Electricity Data Methodology, version 1.6. Kazakhstan source notes. \newline\url{https://files.ember-energy.org/public-downloads/ember_electricity_data_methodology.pdf}

\bibitem{ref3} Samruk-Energy. Electricity market reports, January and January-March 2018. \newline\url{https://samruk-energy.kz/en/press-center/analytical-review}

\bibitem{ref4} Hyndman, R. J. and Athanasopoulos, G. (2021). Forecasting: Principles and Practice, 3rd ed. OTexts. Chapters 5 and 9. \newline\url{https://otexts.com/fpp3/}

\bibitem{ref5} statsmodels. SARIMAX documentation, version 0.15.0. \newline\url{https://www.statsmodels.org/stable/generated/statsmodels.tsa.statespace.sarimax.SARIMAX.html}


\end{thebibliography}

\end{document}
